% SOURCE_STATEMENT: "Does the weak Verlinde formula hold beyond rational conformal field theory; specifically, do asymptotic quantum dimensions define a one-dimensional representation of the fusion ring for nonrational or logarithmic conformal field theories?"
% Preparation and final self-audit: CONTRIBUTING.md prompts C and D; source/audit artifacts retained in the local physics_statements_2026_10_11 batch.
% !TeX program = lualatex
\documentclass[11pt]{article}
\usepackage[margin=25mm]{geometry}
\usepackage{fontspec}
\setmainfont{Latin Modern Roman}
\usepackage{amsmath,amssymb,mathtools}
\usepackage{unicode-math}
\setmathfont{Latin Modern Math}
\usepackage{xurl,hyperref}
\hypersetup{colorlinks=true,urlcolor=blue}
\setcounter{secnumdepth}{0}
\setlength{\parindent}{0pt}
\setlength{\parskip}{5pt}
\setlength{\emergencystretch}{3em}
\newcommand{\sref}[2]{\hyperlink{source:#1:#2}{[S#2]}}
\newcommand{\cataloguescope}[1]{\paragraph{Recorded target.}#1}
\newcommand{\R}{\mathbb R}
\newcommand{\C}{\mathbb C}
\newcommand{\N}{\mathbb N}
\newcommand{\Z}{\mathbb Z}
\newcommand{\Q}{\mathbb Q}
\newcommand{\F}{\mathbb F}
\newcommand{\eps}{\varepsilon}
\begin{document}
\section{Weak Verlinde Formula beyond Rational Conformal Field Theory}\label{statement}

\subsection{Definitions and mathematical statement}

\paragraph{Shared notation.}$\mathbb N=\{1,2,\ldots\}$; the standard number fields have their usual meanings. All additional conventions are stated locally.


\paragraph{Definitions and conventions.}
The fusion ring has a basis of sectors or modules and product given by fusion multiplicities. A one-dimensional representation of this ring is a multiplicative scalar assignment to those basis elements. An asymptotic quantum dimension is obtained from an appropriate limiting ratio of characters, when that limit exists. The weak Verlinde question asks whether these asymptotic dimensions respect fusion even outside semisimple rational conformal theories.

\paragraph{Target retained from the canonical source catalogue.}
\paragraph{Catalogue wording retained for traceability.}
Does the weak Verlinde formula hold beyond rational conformal field theory; specifically, do asymptotic quantum dimensions define a one-dimensional representation of the fusion ring for nonrational or logarithmic conformal field theories?

\paragraph{Source scope and qualification.}
Existence and normalization of the character limits and the chosen fusion ring are part of the problem. Logarithmic and nonrational theories must not be replaced by the rational case in a proposed answer.
\subsection{Short English statement}
Beyond rational conformal theory, do asymptotic quantum dimensions still multiply according to fusion?
\subsection{Sources}
\begin{itemize}
\item[\hypertarget{source:statement:1}{S1}] ULAM.AI and the UnsolvedMath Contributors, UnsolvedMath, catalogue entry. Catalogue curation under CC BY 4.0; definitions and English exposition added locally.
\url{https://huggingface.co/datasets/ulamai/UnsolvedMath}.
\item[\hypertarget{source:statement:2}{S2}] Subfactors and Conformal Field Theory (2015), Oberwolfach Reports, DOI 10.4171/owr/2015/16, pp. 849--926. Primary report accessed 2026-10-11; relevant PDF pages 53, 54.
\url{https://doi.org/10.4171/owr/2015/16}.
\item[\hypertarget{source:statement:3}{S3}] Huang, Verlinde conjecture (catalogue-listed context; not independently consulted in this preparation).
\url{https://arxiv.org/abs/math/0406291}.
\item[\hypertarget{source:statement:4}{S4}] Creutzig--Milas--Wood, regularized dimensions (catalogue-listed context; not independently consulted in this preparation).
\url{https://arxiv.org/abs/1411.3282}.
\end{itemize}
\label{end:statement}
\end{document}
